% =====================================================================
%  Survey: Famous Painting Dynamic Video-Style Generation Techniques
%  Template: IJCAI (International Joint Conference on Artificial Intelligence)
%
%  Compile: pdflatex main.tex && bibtex main && pdflatex main.tex (x2)
%  Note: This file uses the standard IJCAI style file ijcai.sty.
%        Place ijcai.sty in the same folder before compiling.
% =====================================================================
\documentclass[letterpaper,twocolumn,10pt]{article}
\usepackage{ijcai2023} % IJCAI official style; compatible with 2021/2022/2023

\usepackage{times}
\usepackage{helvet}
\usepackage{courier}
\usepackage{url}
\usepackage{hyperref}
\usepackage{graphicx}
\usepackage{amsmath,amssymb}
\usepackage{booktabs}
\usepackage{xcolor}

\title{A Survey on Dynamic Video-Style Generation Techniques for World-Famous Paintings:
From Oil and Ink Style Transfer to Painting Animation}

\author{
  Anonymous Authors\\
  AIforEducation Project\\
  \texttt{aiforeducation@example.com}
}

\begin{document}

\maketitle

\begin{abstract}
World-famous paintings, ranging from Western oil masterpieces of Van Gogh and Rembrandt
to Eastern ink wash landscapes, constitute a unique reservoir of human cultural heritage.
In recent years, the rapid development of deep generative models has opened a new direction:
\emph{dynamic video-style generation for famous paintings}, which aims to bring static
artworks to life through artistic style transfer, image-to-video synthesis, and
portrait/painting animation. This survey provides a structured overview of this emerging
area along four axes: (i) foundations of artistic style transfer; (ii) video style
transfer with temporal consistency; (iii) diffusion-based image and video stylization;
and (iv) painting synthesis, stroke rendering, and ``talking painting'' animation.
We focus on representative works published in top-tier venues including CVPR, ICCV,
ECCV, AAAI, and IJCAI, and discuss open challenges such as content leakage, long-video
coherence, and art-specific aesthetic preservation.
\end{abstract}

% =====================================================================
\section{Introduction}
\label{sec:intro}

Famous paintings encode not only colors and brushstrokes but also the cultural and
emotional intent of their creators. Bringing such static artworks to life---in the form
of short videos---has become a compelling research and educational objective, enabling
museum experiences, AI-driven storytelling, and creative content production. Recent
``living portrait'' demonstrations, in which the \textit{Mona Lisa} or Rembrandt
self-portraits are animated from a single image, illustrate the public appeal of this
direction~\cite{zakharov2019fewshot,guo2024liveportrait}.

Three technological threads have converged to make famous-painting video generation
possible: \textbf{(1) artistic style transfer}, which maps the visual appearance of a
reference painting onto a target image or video; \textbf{(2) image-to-video (I2V) and
video-to-video (V2V) generation}, which animate static content with temporally coherent
motion; and \textbf{(3) painting synthesis and stroke rendering}, which emulate the
physical creation process of paintings, including brushstroke dynamics. This survey
organizes representative works along these axes, with an emphasis on results published at
CVPR, ICCV, ECCV, AAAI, and IJCAI over the past decade.

% =====================================================================
\section{Foundations of Artistic Style Transfer}
\label{sec:style}

\subsection{Optimization-Based and Feed-Forward Methods}

The seminal work of \citet{gatys2016image} demonstrated that the style of an artwork can
be separated from its content using Gram-matrix statistics of deep CNN features, giving
rise to Neural Style Transfer (NST). \citet{johnson2016perceptual} subsequently proposed
a feed-forward network trained with perceptual losses, enabling real-time stylization.
To support arbitrary styles at inference time, \citet{huang2017arbitrary} introduced
Adaptive Instance Normalization (AdaIN), which aligns the channel-wise mean and variance
of content features to those of style features. \citet{park2019arbitrary} further
proposed style-attentional networks (SANet) to attend to relevant style patches.

\subsection{Transformer-Based Arbitrary Style Transfer}

Vision transformers have enabled long-range style modeling.
\citet{deng2021stytr2} proposed StyTr$^2$, which uses a transformer encoder--decoder to
transfer style across arbitrary contents, alleviating the loss of detail common to
AdaIN-based methods. These foundations are the bedrock upon which painting-specific and
video-specific methods are built.

% =====================================================================
\section{Video Style Transfer and Temporal Consistency}
\label{sec:video}

A naive application of image style transfer frame-by-frame produces severe flicker.
Addressing this, \citet{ruder2016artistic} introduced a temporal loss based on optical
flow to constrain consecutive frames, while \citet{chen2017coherent} proposed an online
recurrent framework that propagates stylized features across frames.
\citet{huang2017realtime} presented a real-time feed-forward model for video stylization.

\subsection{Regularization-Based Approaches}

\citet{wang2020consistent} reinterpreted temporal consistency as a compound
regularization problem and designed a zero-shot video style transfer framework,
published at AAAI~2020. \citet{deng2021arbitrary} proposed multi-channel correlation for
arbitrary video style transfer (AAAI~2021), enabling a single network to handle diverse
artistic references without retraining. \citet{gao2020reconet} introduced ReCoNet,
optimizing temporal stability through flow-based warping losses.

\subsection{Diffusion-Based Video Stylization}

The diffusion era has dramatically reshaped video stylization. \citet{ceylan2023pix2video}
and \citet{chai2023stablevideo} (ICCV~2023) showed that image diffusion models can be
extended to video editing by self-attention injection and consistency-aware sampling.
\citet{geyer2023tokenflow} enforced feature consistency across frames through token-level
propagation. More recently, \citet{ye2025stylemaster} proposed StyleMaster (CVPR~2025),
which combines prompt-patch similarity filtering, contrastive style learning, and a
lightweight motion adapter to bridge the image-to-video gap. EchoStyle~\cite{li2026echostyle}
constructs a large-scale stylization dataset of 20k video pairs via reverse synthesis and
achieves performance comparable to commercial closed-source solutions.

% =====================================================================
\section{Diffusion Models for Image and Video Generation}
\label{sec:diffusion}

\subsection{Foundations}

Latent Diffusion Models (LDM)~\cite{rombach2022ldm} established the dominant paradigm
for high-resolution text-to-image synthesis. Controllability was extended by
ControlNet~\cite{zhang2023controlnet}, while subject personalization was enabled by
DreamBooth~\cite{ruiz2023dreambooth} and Textual Inversion~\cite{gal2023textualinversion},
with LoRA~\cite{hu2022lora} providing efficient fine-tuning. These tools form the
practical foundation for painting-style generation pipelines that combine pretrained
diffusion backbones with style LoRAs and ControlNet conditioning.

\subsection{Inversion-Based Style Transfer}

\citet{zhang2023inversion} proposed InST (CVPR~2023), which learns artistic style directly
from a single painting by treating style as a learnable textual description, then guiding
synthesis without complex prompts. ArtAdapter~\cite{chen2024artadapter} introduced a
multi-level style encoder with an explicit adaptation mechanism to capture both low-level
textures and high-level composition.

\subsection{Video Diffusion Models}

Video diffusion models~\cite{ho2022video,blattmann2023svd} extend image diffusion to the
temporal axis. Tune-A-Video~\cite{wu2023tuneavideo} performs one-shot tuning of image
diffusion for text-to-video generation. CoDeF~\cite{ouyang2024codef} and
FlowVid~\cite{liang2024flowvid} (CVPR~2024) introduce content deformation fields and
flow-guided synthesis respectively to ensure temporal coherence. The tuning-free
framework RefVFX~\cite{jones2026refvfx} transfers complex temporal effects from a
reference video, supporting diverse artistic transformations.

% =====================================================================
\section{Painting Synthesis and Stroke Rendering}
\label{sec:painting}

Beyond pixel-level stylization, a complementary line of work models the physical creation
process of paintings.

\subsection{Time-Lapse and Stroke-Level Synthesis}

\citet{zhao2020painting} introduced the task of synthesizing time-lapse videos of
paintings (CVPR~2020), recurrently predicting painting steps from a single completed
artwork. Stylized Neural Painting~\cite{zou2021stylized} (CVPR~2021) formulates
image-to-painting translation as a search for physically meaningful stroke parameters
through a differentiable neural renderer, producing stroke-by-stroke animations.
Neural Strokes~\cite{liu2021neuralstrokes} extends the idea to line drawing of 3D shapes.

\subsection{Data-Driven Brushstroke Rendering}

\citet{madono2023data} proposed a data-driven ink painting brushstroke renderer trained
on a robotically captured dataset of real ink paintings, faithfully reproducing complex
effects such as bleeding and color mixing. Interactive video stylization via few-shot
patch-based training~\cite{texler2020interactive} supports user-guided stylization while
preserving temporal coherence.

% =====================================================================
\section{Art-Style-Specific Methods: Oil and Ink Wash Painting}
\label{sec:artstyle}

\subsection{Oil Painting Style Transfer}

Oil painting is characterized by vivid textures, brush directionality, and rich color
layers. General style transfer methods often produce blurred details and texture loss on
oil paintings, motivating specialized architectures that incorporate stroke-texture
constraints, adaptive content--style weighting, and color statistical alignment.
Combined with diffusion backbones and painter-specific LoRAs (e.g., Van Gogh, Monet,
Rembrandt), modern pipelines achieve substantially higher artistic fidelity.

\subsection{Chinese Ink Wash and Shanshui Painting}

Chinese ink wash painting poses unique challenges: it follows the classical five-level
ink theory, employs strategic negative space, and uses brushstrokes that convey both
form and spirit. Instance-Aware Coherent Video Style Transfer for Chinese Ink Wash
Painting~\cite{liang2021instance} (IJCAI~2021) addresses object-level consistency in
ink-wash video stylization. BiTGAN~\cite{he2023bitgan} combines a UNet path and a ResNet
path in a bilateral generator to balance structural preservation and ink texture.
CVSTGAN~\cite{wang2024cvstgan} introduces a controllable GAN for Chinese painting video
style transfer. ChipDiff~\cite{liu2026chipdiff} applies a staged diffusion model with
loss-gradient guidance to better capture ink-specific aesthetics. InkSpirit~\cite{zhang2025inkspirit}
incorporates expert knowledge---blank space principles and inscription semantics---into a
ComfyUI-based text-to-image framework for traditional Chinese painting.

% =====================================================================
\section{Portrait and ``Talking Painting'' Animation}
\label{sec:portrait}

A particularly compelling application is animating famous painted portraits into
``talking paintings.'' \citet{zakharov2019fewshot} introduced few-shot adversarial
learning of neural talking head models, enabling animations of the \textit{Mona Lisa}
and other portraits from a single image. LivePortrait~\cite{guo2024liveportrait}
achieves efficient portrait animation through implicit keypoints with stitching and
retargeting control, generalizing to stylized portraits such as oil paintings and anime.
EDTalk~\cite{tan2024edtalk} (ECCV~2024) disentangles mouth shape, head pose, and emotion
for fine-grained control. The MEAD dataset~\cite{wang2020mead} (ECCV~2020) provides a
large-scale emotional talking-face benchmark that supports training emotion-aware
animated portraits of artworks.

% =====================================================================
\section{Discussion and Open Challenges}
\label{sec:discussion}

\textbf{Content leakage and style--content disentachment.} Image-trained stylization
often transfers semantic content from the style reference, distorting the target.
Recent works~\cite{ye2025stylemaster,chen2024artadapter} tackle this through patch
filtering and auxiliary content adapters, but the problem remains open for long videos.

\textbf{Long-video temporal coherence.} Sliding-window strategies~\cite{li2026echostyle}
and content deformation fields~\cite{ouyang2024codef} mitigate drift but struggle with
large motion and complex scenes. Cross-turn consistency in iterative editing is a
related open problem.

\textbf{Art-specific aesthetic preservation.} General style transfer fails to capture
art-specific features such as ink density, brush direction, and compositional blank
space. Specialized losses and datasets~\cite{he2023bitgan,liu2026chipdiff,zhang2025inkspirit}
are essential but currently limited to specific art traditions.

\textbf{Integration of generation and animation pipelines.} Combining painting-specific
style LoRAs, ControlNet conditioning, I2V diffusion models, and talking-portrait
animation into a single pipeline---as exemplified by the OilPaintingVideo project---is
an emerging engineering direction with significant educational potential.

% =====================================================================
\section{Conclusion}
\label{sec:conclusion}

Dynamic video-style generation for world-famous paintings is a fast-evolving intersection
of artistic style transfer, diffusion-based video generation, stroke-level painting
synthesis, and portrait animation. This survey has reviewed representative works across
CVPR, ICCV, ECCV, AAAI, and IJCAI, organizing them along four axes. We believe that the
convergence of pretrained diffusion backbones with art-specific knowledge and temporal
consistency mechanisms will continue to drive progress in this culturally significant
direction, with broad applications in education, museums, and creative industries.

% =====================================================================
\bibliographystyle{named}
\bibliography{references}

\end{document}
